% Three Smoothing Mechanisms in Early Cosmology
% LuaLaTeX / XeLaTeX with TeX Gyre Pagella
% Author: Flyxion (Independent Researcher)

\documentclass[12pt, a4paper]{article}
\usepackage[T1]{fontenc}
\usepackage{mathpazo}


\usepackage{microtype}
\usepackage{amsmath, amssymb, amsthm}
\usepackage{mathtools}
\usepackage{geometry}
\geometry{margin=1.25in}
\usepackage{setspace}
\onehalfspacing
\usepackage{hyperref}
\hypersetup{colorlinks=true, linkcolor=black, citecolor=black, urlcolor=black}
\usepackage{natbib}
\bibpunct{[}{]}{,}{n}{}{;}
\usepackage{booktabs}
\usepackage{array}

% Theorem environments
\newtheorem{definition}{Definition}
\newtheorem{proposition}{Proposition}
\newtheorem{remark}{Remark}

% Custom commands
\newcommand{\vphi}{\phi}
\newcommand{\vv}{\mathbf{v}}
\newcommand{\calA}{\mathcal{A}}
\newcommand{\calF}{\mathcal{F}}
\newcommand{\etalamph}{\eta_{\rm lamph}}
\newcommand{\epslamph}{\epsilon_{\rm lamph}}
\newcommand{\Heff}{H_{\rm eff}}
\newcommand{\rRSVP}{r_{\rm RSVP}}
\newcommand{\rQQG}{r_{\rm QQG}}
\newcommand{\Ph}{P_h}
\newcommand{\Vol}{\operatorname{Vol}}
\newcommand{\ddt}{\frac{d}{dt}}
\newcommand{\Teff}{T^{\rm eff}_{ij}}

\title{\textbf{Three Smoothing Mechanisms in Early Cosmology:\\
Inflaton, Geometric, and Lamphrodyne Relaxation}}

\author{Flyxion\\[4pt]
\small Independent Researcher}

\date{June 2026}

\begin{document}

\maketitle

\begin{abstract}
We propose a classification of early-universe smoothing mechanisms into three
structurally distinct classes: inflaton smoothing, in which an auxiliary scalar
field with tuned potential drives metric expansion; geometric smoothing, in which
quantum-corrected gravitational dynamics generate an expansion phase without a
separate inflaton; and lamphrodyne relaxation, in which cosmological homogenization
arises from admissibility-reducing dynamics in a scalar-vector-entropy plenum,
without global metric expansion. The recent result of Liu, Quintin, and Afshordi---
quantum quadratic gravity (QQG) inflation published in \textit{Physical Review
Letters} 136, 111501 (2026)---establishes geometric smoothing as a theoretically
coherent alternative to the inflaton paradigm. We treat QQG as a rigorous benchmark
and ask the next question: if smoothing can emerge from quantum gravity alone, can
it emerge from relaxation dynamics alone? We formalize the Relativistic
Scalar-Vector Plenum (RSVP) framework as a candidate Class~3 mechanism, define the
lamphrodyne background state, introduce the field-space admissibility metric, and
derive the lamphrodyne relaxation parameter $\etalamph$ from admissibility geometry.
We construct the effective anisotropic stress tensor of the relaxing plenum and
define the formal structure of the tensor power spectrum $\Ph(k)$ and
tensor-to-scalar ratio $\rRSVP$; their explicit evaluation requires the plenum
field correlators, which the prototype action makes calculable in principle but
which are not derived here. We provide a phenomenological estimate of the scalar
spectral index using a multifield GR analogy, and introduce admissibility crossing
as the RSVP-native freeze-out condition that a complete perturbation theory must
employ. We identify an empirical fork: a detection of $r > 10^{-2}$ constrains
RSVP to explain any tensor residue from plenum anisotropic stress, while a robust
upper bound below this floor would strongly challenge the simplest QQG realization.
The paper's primary contribution is not adjudicating between frameworks but
establishing that Class~3 smoothing is formally tractable as an effective field
theory program with defined observational targets and explicit falsification
conditions.
\end{abstract}

\tableofcontents
\newpage

%------------------------------------------------------------
\section{Introduction}
\label{sec:intro}
%------------------------------------------------------------

\subsection{The Inflationary Paradigm}

The standard cosmological model inherits three observational requirements that are
not explained by the hot big bang alone. The horizon problem asks why causally
disconnected regions of the cosmic microwave background (CMB) share a temperature
to one part in $10^5$. The flatness problem asks why the spatial curvature of the
universe is tuned to within observational bounds across many orders of magnitude in
scale factor. The structure problem asks why primordial density perturbations exhibit
a nearly scale-invariant power spectrum. Inflation---a phase of accelerated expansion
in the very early universe---was introduced to resolve all three simultaneously
\citep{Guth1981, Linde1982, Albrecht1982}.

In its standard formulation, inflation is driven by a scalar field $\phi$, the
inflaton, rolling slowly down a potential $V(\phi)$. The slow-roll conditions
$\epsilon \ll 1$ and $|\eta| \ll 1$, where
\begin{equation}
\epsilon = \frac{M_{\rm Pl}^2}{2}\left(\frac{V'}{V}\right)^2, \qquad
\eta = M_{\rm Pl}^2 \frac{V''}{V},
\end{equation}
ensure that the potential energy dominates over kinetic energy, sustaining a
quasi-de Sitter expansion phase. The inflationary paradigm is empirically powerful:
it predicts a nearly scale-invariant scalar spectral index $n_s \approx 1 - 2/N$
and a tensor-to-scalar ratio $r \approx 16\epsilon$, both in agreement with CMB
observations at leading order \citep{Planck2020}.

\subsection{The Inflaton Problem}

Despite its empirical success, the inflaton paradigm carries well-known theoretical
liabilities. The shape of $V(\phi)$ must be specified by hand; no derivation from a
UV-complete theory of gravity selects a unique potential. The inflationary
literature contains hundreds of viable potential forms, most differing only in their
predictions for $r$ and $n_s$ at sub-percent precision. This proliferation---the
landscape of inflationary models---is not a scientific failure so much as a symptom
of under-determination: the paradigm is consistent with almost any slow-roll
potential that fits a plateau \citep{Martin2014}.

Reheating presents a second ambiguity. Inflation must end, and the universe must
transition from a cold, potential-energy-dominated state to the hot radiation-
dominated era of standard big-bang nucleosynthesis. The mechanism of this
transition---perturbative decay of the inflaton, parametric resonance, instant
preheating---is not predicted by the inflationary framework itself but must be
appended \citep{Kofman1997}. The reheating temperature $T_{\rm reh}$ is largely
unconstrained, introducing a significant uncertainty in the mapping between
inflationary observables and CMB scales.

\subsection{The QQG Development}

The recent result of \citet{Liu2026} represents a qualitatively different approach.
Rather than introducing an inflaton field, they begin with the action of quadratic
gravity,
\begin{equation}
S_{\rm QQG} = -\int d^4x\,\sqrt{-g}\left(\frac{R^2}{\xi} + \frac{C^2}{2\lambda}\right),
\end{equation}
where $R$ is the Ricci scalar, $C^2 \equiv C_{\mu\nu\rho\sigma}C^{\mu\nu\rho\sigma}$
is the Weyl tensor contraction, and $\xi$, $\lambda$ are dimensionless coupling
constants. This theory is renormalizable---in contrast to general relativity---and
is asymptotically free in the ultraviolet (UV), meaning the couplings vanish as
energy $\mu \to \infty$ \citep{Fradkin1982, Avramidi1985}.

The inflationary phase in QQG does not arise from a classical potential but from
one-loop renormalization group (RG) running of the coupling $\xi(\mu)$. Pure $R^2$
gravity is exactly scale-invariant; any constant-curvature spacetime---including de
Sitter---extremizes the action. The running breaks this exact invariance through a
quantum conformal anomaly \citep{Capper1974}, generating a logarithmic correction
that stabilizes a quasi-de Sitter attractor. In the large-$N$ limit, with $N \sim
\mathcal{O}(10^5$--$10^6)$ matter fields contributing to the beta functions, the
spectral index and tensor-to-scalar ratio are predicted to lie in the region
preferred by current CMB data, including the ACT results that place Starobinsky
inflation under mild tension \citep{Louis2025}.

A structural prediction of the QQG scenario is a minimum tensor-to-scalar ratio
\begin{equation}
\rQQG \gtrsim 10^{-2},
\end{equation}
arising from the requirement that the theory avoid its strong-coupling regime before
reheating is complete. This bound is not a free parameter but a consequence of the
theory's internal consistency conditions.

\subsection{A Different Question}

The significance of QQG for the present paper is not its viability as an
inflationary model. Its significance is structural: QQG demonstrates that
cosmological smoothing can emerge from the intrinsic dynamics of a fundamental
theory, without an auxiliary inflaton field. Once this possibility is admitted,
inflationary expansion becomes one member of a broader class of smoothing
mechanisms rather than the unique solution to the horizon and flatness problems.

This paper asks whether a third class of smoothing mechanism is formally tractable:
relaxation without expansion. The RSVP framework---based on a coupled
scalar-vector-entropy plenum---proposes that early-universe homogenization arises
from the directional reduction of admissible unresolved field configurations, driven
by entropy production in the plenum rather than by metric expansion. We formalize
this proposal, derive its analogue of slow-roll parameters, and identify the
observational signatures that distinguish it from both Class~1 (inflaton) and
Class~2 (geometric) smoothing.

%------------------------------------------------------------
\section{A Taxonomy of Smoothing Mechanisms}
\label{sec:taxonomy}
%------------------------------------------------------------

\subsection{What Requires Explanation}

The observational targets that any smoothing mechanism must address are the
following. First, CMB isotropy: the temperature anisotropy $\Delta T / T \sim
10^{-5}$ across the last scattering surface, including regions that were not in
causal contact under standard big-bang evolution. Second, large-scale homogeneity:
the matter distribution is statistically isotropic and homogeneous on scales above
approximately $100\,h^{-1}$~Mpc. Third, the primordial power spectrum: scalar
perturbations are nearly scale-invariant with spectral index $n_s \approx 0.965$
and amplitude $A_s \approx 2.1 \times 10^{-9}$ \citep{Planck2020}. Fourth,
primordial tensor modes: no detection as of 2026, with current upper bounds
$r < 0.036$ at 95\% confidence \citep{BICEPKeck2021}.

\subsection{Three Classes of Solution}

\textit{Class 1: Inflaton smoothing.} An auxiliary scalar field $\phi$ with
potential $V(\phi)$ drives a phase of accelerated metric expansion. The expansion
causally connects regions that would otherwise be disconnected. Smoothing is
achieved by dilution: any initial anisotropy or curvature is stretched to scales
far beyond the Hubble horizon. The slow-roll parameters $\epsilon$ and $\eta$
control the duration and observables of the inflationary phase.

\textit{Class 2: Geometric smoothing.} No auxiliary scalar is introduced. An
inflationary phase emerges from quantum corrections to the gravitational action
itself. In QQG, the one-loop RG running of the coupling $\xi(\mu)$ breaks the
exact scale invariance of $R^2$ gravity, generating a quasi-de Sitter attractor.
The mechanism is geometrical in the sense that no matter field carries the
inflationary energy; the gravitational sector alone generates the expansion.
Nevertheless, the output is still a period of accelerated metric expansion, and
the observational predictions---scalar spectrum, tensor spectrum, spectral
tilt---are structurally similar to Class~1, differing in their specific values of
$n_s$ and $r$.

\textit{Class 3: Lamphrodyne relaxation.} No metric expansion is invoked.
Cosmological smoothing arises from the directional relaxation of a coupled
scalar-vector-entropy plenum toward states of lower admissibility complexity. The
apparent homogeneity of the CMB is not the result of causal contact established
during an expansion phase but the thermodynamic equilibrium texture of a plenum
that has undergone irreversible entropy-producing relaxation. Cosmological redshift
is reinterpreted as an entropy-gradient effect in the scalar field $S$ rather than
a consequence of metric expansion.

\subsection{Why Three Classes?}

Any smoothing mechanism must answer two fundamental questions: what drives
homogenization, and is metric expansion required to achieve it? These two questions
generate the classification naturally. If homogenization requires expansion, the
question becomes whether the expansion is driven by an auxiliary matter field
(Class~1) or by the gravitational sector itself (Class~2). If homogenization does
not require expansion, a third possibility opens (Class~3), in which field-space
relaxation produces the same observational regularities without a de Sitter phase.

\begin{center}
\begin{tabular}{lll}
\toprule
\textbf{Class} & \textbf{Driver} & \textbf{Expansion required?} \\
\midrule
1 (Inflaton) & Auxiliary scalar $\phi$ & Yes \\
2 (Geometric) & Quantum-corrected gravity & Yes \\
3 (Lamphrodyne) & Admissibility relaxation & No \\
\bottomrule
\end{tabular}
\end{center}

The classification is not ad hoc. It follows from asking which of the two questions
is answered differently. Class~1 and Class~2 share the same answer to the expansion
question and differ only in the source of the driving energy. Class~3 gives a
different answer to both questions. It is therefore a genuinely distinct mechanism,
not a variant of inflation with a non-standard energy source.

\subsection{Ontology Versus Mechanism}

The distinction between Class~2 and Class~3 is not merely technical. Both replace
the inflaton with something more fundamental; they disagree about whether
``fundamental'' means a quantum-corrected metric or a relaxing plenum. Class~2
retains metric expansion as the physical mechanism of smoothing, deriving the
driving force from geometry rather than matter. Class~3 rejects the necessity of
expansion altogether, proposing that the observational signatures attributed to
inflation---isotropy, homogeneity, scale-invariant perturbations---can be generated
by admissibility-reducing dynamics in field space without the universe ever
undergoing a de Sitter phase.

This distinction carries empirical content. If the early universe underwent a
de Sitter phase, whether driven by an inflaton (Class~1) or by quantum gravity
(Class~2), primordial tensor modes are generated by quantum fluctuations of the
metric during that phase. If the early universe instead underwent lamphrodyne
relaxation (Class~3), tensor modes arise from a different source---anisotropic
stress in the relaxing plenum---and their amplitude, tilt, and spectral shape may
differ from the inflationary prediction.

%------------------------------------------------------------
\section{Quadratic Quantum Gravity as Benchmark}
\label{sec:qqg}
%------------------------------------------------------------

Class~2 geometric smoothing encompasses any framework in which an inflationary
phase emerges from the quantum-corrected gravitational sector without an auxiliary
inflaton. Several such proposals exist, including Starobinsky $R + R^2$ inflation
\citep{Starobinsky1980}, asymptotic safety scenarios, induced gravity models, and
higher-curvature extensions of GR. We use the specific QQG scenario of
\citet{Liu2026} as our primary benchmark because it is the most recent Class~2
model to make a structural falsification prediction---the tensor floor $r \gtrsim
10^{-2}$---and because it has been tested against the most current CMB data
including ACT DR6. The classification of smoothing mechanisms developed here is
independent of the fate of this particular model; if the Liu--Quintin--Afshordi
scenario were falsified, the Class~2 category would remain and would be occupied
by other geometric smoothing proposals. The empirical fork described in
Section~\ref{sec:fork} should therefore be read as a comparison between Class~2
and Class~3 smoothing in general, with QQG serving as the sharpest available
Class~2 representative.

\subsection{Action and Scale Invariance}

The QQG action of \citet{Liu2026} is equation~(1) above. In homogeneous and
isotropic backgrounds, the Weyl tensor $C_{\mu\nu\rho\sigma}$ vanishes identically,
so the background dynamics are controlled entirely by the $R^2/\xi$ term. Without
RG running, pure $R^2$ gravity is exactly scale invariant: both Minkowski and
(anti-)de Sitter spacetimes are solutions. This scale invariance is analogous to
the classical conformal invariance of Yang-Mills theory before the introduction of
a mass scale through dimensional transmutation.

\subsection{RG Flow and Dimensional Transmutation}

The one-loop beta functions derived by \citet{Buccio2024} are
\begin{align}
\beta_\xi &= \frac{d\xi}{d\ln\mu} = -\frac{1}{(4\pi)^2}\frac{\xi^2 - 36\lambda\xi - 2520\lambda^2}{36}, \\
\beta_\lambda &= \frac{d\lambda}{d\ln\mu} = -\frac{1}{(4\pi)^2}\frac{(1617 + 90N)\lambda - 20\xi\lambda}{90},
\end{align}
where $N$ encodes the matter field content through $N = \frac{1}{60}N_{\rm scalar} +
\frac{1}{5}N_{\rm vector} + \frac{1}{20}N_{\rm fermion}$. The running coupling
$\xi(\mu)$ exhibits a maximum at some scale $\mu_{\rm max}$ and subsequently
decreases, crossing zero at the tachyon divide $\mu_0$. \citet{Liu2026} identify
the inflationary epoch with the regime of decreasing $\xi$ above the tachyon divide.

In the large-$N$ limit, the solution is approximately
\begin{equation}
\xi(\mu) \simeq \frac{35\lambda_0^2\ln(\mu/\mu_0)}{8\pi^2[1 + \lambda_{tH}\ln(\mu/\mu_0)]},
\end{equation}
where $\lambda_{tH} \equiv \lambda_0 N/(4\pi)^2$ is a `t~Hooft-like coupling. The
inflation potential in the Einstein frame takes the approximate form
\begin{equation}
V(\varphi) \simeq \frac{35\lambda_0^2\mu_0^4}{128\pi^2\lambda_{tH}}
\left(1 - \frac{\sqrt{6}\,\mu_0}{\lambda_{tH}\varphi}\right),
\end{equation}
a slowly varying plateau that supports slow-roll inflation without a separately
introduced inflaton field.

\subsection{Strong Coupling and Ghost Confinement}

The Weyl term $C^2$ introduces a massive spin-2 ghost: a field carrying negative
kinetic energy. In perturbative QQG this field generates an unbounded Hamiltonian
if treated in isolation. \citet{Liu2026} adopt the proposal of \citet{Holdom2016}
that, analogously to quarks in QCD, ghost degrees of freedom are confined once QQG
reaches its strong-coupling scale. The end of inflation, the tachyon divide, and
the onset of strong coupling are argued to be nearly coincident---a result that is
empirically suggested by the parameter analysis but remains a conjecture awaiting
nonperturbative verification.

\subsection{Emergence of General Relativity}

Perhaps the most radical claim in \citet{Liu2026} is not that inflation can be
derived from quantum gravity. It is that general relativity itself is not
fundamental. The QQG framework proposes that Einstein's theory emerges as an
infrared effective field theory from the strong-coupling limit of a UV-complete
renormalizable theory. The logical chain is:
\begin{equation}
\text{UV Fixed Point} \;\rightarrow\; \text{RG Running} \;\rightarrow\;
\text{Quasi-de Sitter Inflation} \;\rightarrow\; \text{Strong Coupling} \;\rightarrow\;
\text{Emergent GR}.
\end{equation}
This is conceptually analogous to the emergence of nuclear physics from QCD: the
low-energy effective description (GR, nuclear physics) is not a fundamental theory
but a phase of a deeper one. In QQG, the Planck mass $M_{\rm Pl}$ is not a
fundamental input but an emergent scale arising when the 't~Hooft-like coupling
$\lambda_{tH}$ surpasses unity and the perturbative QQG description breaks down.
An Einstein-Hilbert term with effective $M_{\rm Pl}$ then appears as the leading
term of the low-energy effective action, and the universe enters its familiar
radiation-dominated era.

The title of \citet{Liu2026}---``Ultraviolet Completion of the Big Bang''---reflects
this ambition. The paper is not primarily proposing a new inflation model; it is
proposing a new origin for the entire gravitational framework within which inflation
operates. This framing is directly relevant to RSVP, which shares the structural
intuition that the familiar low-energy description is not fundamental: in QQG,
spacetime geometry is emergent from a UV gravitational phase; in RSVP, spacetime
geometry is an effective description of plenum dynamics. The disagreement between
the frameworks is over what lies beneath, not over whether the familiar description
is fundamental.

\subsection{Observable Predictions}

The scalar spectral index and tensor-to-scalar ratio predicted by QQG are
\begin{equation}
n_s \sim 1 - \frac{4}{3N_e}, \qquad r \sim 8\left(\frac{2\lambda_{tH}^2}{N_e^4}\right)^{1/3},
\end{equation}
where $N_e \sim 50$--$60$ is the number of inflationary e-folds. The viable
parameter space requires $N \sim 10^5$--$10^6$ matter fields and
$\lambda_{tH} \in [0.1, 1]$. The structural lower bound on the tensor-to-scalar
ratio,
\begin{equation}
\rQQG \gtrsim 10^{-2},
\label{eq:qqg_floor}
\end{equation}
follows from $r \sim \lambda_{tH}^{-2/3}$ and the perturbativity bound
$\lambda_{tH} \lesssim 1$. This bound is falsifiable by next-generation B-mode
experiments including LiteBIRD and CMB-S4.

%------------------------------------------------------------
\section{RSVP Cosmology}
\label{sec:rsvp}
%------------------------------------------------------------

\subsection{The Scalar-Vector Plenum}

The Relativistic Scalar-Vector Plenum (RSVP) describes the early universe through
three coupled fields defined on a four-dimensional spacetime manifold $\mathcal{M}$:
a scalar field $\Phi$ encoding local energy density and constraint-satisfaction; a
vector field $\vv$ encoding directed entropy flux and local flow; and a scalar
entropy field $S$ encoding thermodynamic irreversibility and admissibility state.
The field equations governing their joint evolution are derived from an action
principle supplemented by admissibility constraints, the latter encoding the
requirement that the plenum remain in a physically accessible region of field space.

The RSVP framework does not begin with a metric ansatz. Rather, the effective
metric perceived by observers arises from the gradient structure of $S$: regions
in which $\nabla S$ is large correspond to what standard observers would interpret
as high-redshift regions, since photon traversal across an entropy gradient
accumulates a frequency shift without metric expansion. The effective Hubble law in
RSVP is therefore a statement about entropy-gradient accumulation rather than
recession velocity.

\subsection{Redshift Without Expansion}

In standard cosmology, the cosmological redshift of a photon traveling from emission
redshift $z_e$ to an observer is attributed to the stretching of its wavelength by
metric expansion: $1 + z = a(t_{\rm obs})/a(t_{\rm em})$. In RSVP, no global scale
factor $a(t)$ is introduced. Instead, a photon traversing a region of entropy
gradient $\nabla S$ accumulates a phase shift proportional to the integrated entropy
contrast along its path. The effective redshift is
\begin{equation}
1 + z_{\rm RSVP} = \exp\left(\int_{\rm path} \kappa\,|\nabla S|\,d\ell\right),
\end{equation}
where $\kappa$ is a coupling constant between the photon propagator and the entropy
field, and $d\ell$ is the proper path element. At leading order in a slowly varying
$S$, this reproduces a Hubble-like distance-redshift relation with effective Hubble
parameter $\Heff \sim \kappa\langle|\nabla S|\rangle$.

This reinterpretation is not merely definitional. It predicts that the
distance-redshift relation will exhibit small deviations from the standard
$\Lambda$CDM prediction at high redshift, where the entropy gradient structure of
the plenum departs from homogeneity. These deviations are in principle observable
with future spectroscopic surveys, though their amplitude depends on the specific
background solution $S_0(t)$ derived in Section~\ref{sec:background}.

%------------------------------------------------------------
\section{Lamphrodyne Background Dynamics}
\label{sec:background}
%------------------------------------------------------------

\subsection{The Homogeneous Relaxation State}

We define the lamphrodyne epoch as the cosmological phase during which the RSVP
plenum undergoes homogeneous, isotropic, irreversible relaxation from an initial
state of high admissibility complexity toward a state of low admissibility
complexity. The background is characterized by spatially homogeneous field
configurations
\begin{equation}
\Phi = \Phi_0(t), \qquad \vv = \vv_0(t) = \mathbf{0}, \qquad S = S_0(t),
\end{equation}
where the vanishing of $\vv_0$ follows from the isotropy of the background: a
nonzero background vector field would select a preferred spatial direction,
violating statistical isotropy. The vector field $\vv$ therefore contributes to the
background only through its fluctuations, which are treated in
Section~\ref{sec:perturbations}.

The background evolution equations for $\Phi_0(t)$ and $S_0(t)$ follow from the
RSVP field equations in the homogeneous limit. Their specific form depends on the
choice of kinetic terms and coupling functions in the RSVP action, which we leave
partially general here while imposing the constraints required by admissibility
geometry.

\subsection{The Irreversibility Condition}

The defining property of the lamphrodyne epoch is that the volume of the
unresolved admissible configuration space $\calA_{\rm unresolved}$ decreases
monotonically in time:
\begin{equation}
\frac{d}{dt}\Vol(\calA_{\rm unresolved}) < 0.
\label{eq:irreversibility}
\end{equation}
This is the plenum-theoretic statement of the second law of thermodynamics during
the relaxation epoch. It rules out background solutions along which admissibility
volume increases---which would correspond to a physical unsmoothing of the
plenum---and selects the thermodynamically privileged direction of evolution.

Equation~\eqref{eq:irreversibility} provides a natural arrow of time for the
lamphrodyne epoch that is absent from inflationary frameworks. In Class~1 and
Class~2 smoothing, the thermodynamic arrow of time must be imposed separately from
the dynamical equations; inflation itself is time-reversible at the level of the
field equations. In RSVP, the arrow of time is primitive: relaxation is
admissibility loss, and admissibility loss is by definition irreversible.

\subsection{Admissibility Volume and Field-Space Geometry}
\label{sec:fieldspace}

Let $\calF = \{(\Phi, \vv, S)\}$ denote the field-space manifold of the RSVP
plenum. We equip $\calF$ with a Riemannian metric $G_{AB}$ derived from the kinetic
terms of the RSVP action, so that the kinetic energy of the plenum fields takes the
form
\begin{equation}
\mathcal{T} = \frac{1}{2}G_{AB}\dot{\phi}^A\dot{\phi}^B,
\end{equation}
where $\phi^A = (\Phi, v^i, S)$ collectively denotes the plenum field components.
The metric $G_{AB}$ is positive definite when the kinetic terms are non-degenerate,
which we assume throughout.

The admissibility volume functional $\Vol(\calA)$ assigns to each point
$\phi^A \in \calF$ a real number measuring the volume of the set of configurations
consistent with the admissibility constraints at that point. The gradient
\begin{equation}
\nabla_{\calF}\log\Vol(\calA)
\end{equation}
is a vector field on $\calF$ pointing in the direction of steepest increase of log-
admissibility-volume. During lamphrodyne relaxation, the plenum evolves in the
direction of decreasing $\Vol(\calA)$, which is opposite to this gradient.

Concretely, the admissibility volume at a point $\phi^A \in \calF$ is defined by
\begin{equation}
\Vol(\calA) = \int_{\calA} \sqrt{\det G}\; d\Phi\, dS\, d^3v,
\end{equation}
where $\sqrt{\det G}$ is the Riemannian volume element of the field-space metric
$G_{AB}$. This expression is manifestly geometric: it measures the metric volume of
the admissible region in $\calF$, making $\Vol(\calA)$ a coordinate-invariant
quantity rather than a metaphorical one. As the plenum relaxes and constraints are
resolved, the set $\calA$ shrinks and $\Vol(\calA)$ decreases, in accordance with
the irreversibility condition~\eqref{eq:irreversibility}.

To make the geometry concrete, consider the simplest non-degenerate field-space
metric consistent with the RSVP field content:
\begin{equation}
ds^2_{\calF} = d\Phi^2 + \alpha_v\,|d\vv|^2 + \alpha_S\,dS^2,
\label{eq:toy_metric}
\end{equation}
where $\alpha_v > 0$ and $\alpha_S > 0$ are positive coupling constants encoding
the relative contributions of the vector and entropy fields to the field-space
geometry. In this metric, $\det G = \alpha_v^3\,\alpha_S$ (with three spatial
components of $\vv$), and the admissibility volume element is
$\sqrt{\det G}\,d\Phi\,dS\,d^3v = \alpha_v^{3/2}\alpha_S^{1/2}\,d\Phi\,dS\,d^3v$.
The physical RSVP metric $G_{AB}$ will in general depend on the fields themselves
and on the admissibility constraints; equation~\eqref{eq:toy_metric} is a
linearized approximation valid near a background point in $\calF$.

\subsection{Irreversibility and Admissibility Loss}
\label{sec:irreversibility}

In the Spherepop formalization of irreversible event calculus, the collapse operator
$\mathsf{Pop}$ and the refusal operator $\mathsf{Refuse}$ describe the asymmetric
resolution of unresolved configurations: a configuration either collapses to a
resolved state (Pop) or is refused (Refuse), but neither operation is reversible.
At the level of the lamphrodyne background, this structure is captured by the
single condition~\eqref{eq:irreversibility}: the set of unresolved admissible
configurations shrinks. We do not require the full Spherepop apparatus for the
present calculation; we require only that the background evolution satisfies the
irreversibility condition.

The irreversibility condition also implies that the background trajectory in $\calF$
is not time-reversible. This is in contrast to the de Sitter solution in QQG, which
is an exact solution of the classical equations and is time-symmetric. The
lamphrodyne background is thermodynamically directed from the outset, which is the
formal ground for the claim that Class~3 smoothing provides an arrow of time
without additional postulation.

\begin{proposition}[Monotonic Admissibility Loss]
Suppose that $\Vol(\calA_{\rm unresolved})$ is bounded below by zero and that the
irreversibility condition $d\Vol(\calA_{\rm unresolved})/dt < 0$ holds throughout
the lamphrodyne epoch. Then $\Vol(\calA_{\rm unresolved})$ converges to a finite
non-negative limit as $t \to \infty$.
\end{proposition}

\begin{proof}
The function $t \mapsto \Vol(\calA_{\rm unresolved}(t))$ is strictly decreasing by
the irreversibility condition and bounded below by zero by assumption. A monotone
decreasing function bounded below converges. \qed
\end{proof}

\begin{remark}
The limiting value $\Vol(\calA_{\rm unresolved}) \to V_\infty \geq 0$ represents
the irreducible residual admissibility volume: the set of configurations that cannot
be resolved by the lamphrodyne dynamics alone. If $V_\infty = 0$, the relaxation is
complete and the plenum reaches a fully resolved state. If $V_\infty > 0$, a
permanently unresolved residue persists, which may correspond to the observed
large-scale structure of the universe---the non-trivial texture that survives after
smoothing is complete.
\end{remark}

%------------------------------------------------------------
\section{The Lamphrodyne Relaxation Parameter}
\label{sec:etalamph}
%------------------------------------------------------------

\subsection{Motivation by Analogy}

In slow-roll inflation, the parameters $\epsilon$ and $\eta$ measure, respectively,
the steepness and curvature of the inflaton potential. Their smallness is the
condition under which the inflationary background is approximately de Sitter and
perturbation theory is under control. The spectral observables $n_s$ and $r$ are
directly expressed in terms of $\epsilon$ and $\eta$ through the consistency
relations of inflationary perturbation theory.

For RSVP, the analogous role is played by a parameter measuring the rate at which
the admissibility relaxation gradient changes during the lamphrodyne epoch. We
define this parameter from admissibility geometry rather than by importing the
slow-roll definition.

\subsection{Definition}

\begin{definition}[Lamphrodyne Gradient Parameter]
Let $\calF$ be the field-space manifold of the RSVP plenum with metric $G_{AB}$,
and let $\Vol(\calA)$ be the admissibility volume functional. Define the
\emph{lamphrodyne gradient parameter} as
\begin{equation}
\epslamph = \frac{\left\|\nabla_{\calF}\log\Vol(\calA)\right\|}{\Heff},
\label{eq:epslamph}
\end{equation}
where $\Heff \sim \kappa\langle|\nabla S|\rangle$ is the effective Hubble parameter
from entropy-gradient traversal, and the norm is taken with respect to the
field-space metric $G_{AB}$. The parameter $\epslamph$ measures the magnitude of
the admissibility gradient in units of $\Heff$: it is the lamphrodyne analogue of
the slow-roll parameter $\epsilon_{\rm infl}$, measuring the steepness of the
admissibility descent.
\end{definition}

\begin{definition}[Lamphrodyne Relaxation Parameter]
Define the \emph{lamphrodyne relaxation parameter} as
\begin{equation}
\etalamph = -\frac{1}{\Heff}\frac{d}{dt}\log\epslamph,
\label{eq:etalamph}
\end{equation}
measuring the fractional rate of change of $\epslamph$, normalized by $\Heff$.
Small $|\etalamph|$ implies that the admissibility gradient is nearly constant in
time---the lamphrodyne analogue of the slow-roll condition $|\eta_{\rm infl}| \ll 1$.
\end{definition}

\subsection{Physical Interpretation}

The primary new object introduced in this framework is $\epslamph$. It measures the
magnitude of the admissibility gradient in field space, normalized by the effective
Hubble parameter: it is the generalized descent rate of the relaxation trajectory on
$\calF$. In this respect $\epslamph$ plays exactly the role that $\epsilon_{\rm infl}$
plays in slow-roll inflation: it quantifies how steeply the system is descending its
governing landscape, whether that landscape is a scalar potential (inflation) or an
admissibility volume functional (lamphrodyne relaxation).

The parameter $\etalamph$ is then a derived curvature measure: it quantifies how
rapidly $\epslamph$ itself is changing. Small $|\etalamph|$ means that the descent
rate is nearly constant---the relaxation is in a controlled, nearly uniform regime.
This is the lamphrodyne analogue of the slow-roll condition $|\eta_{\rm infl}| \ll 1$,
which ensures that $\epsilon_{\rm infl}$ varies slowly enough for inflationary
perturbation theory to be valid.

The hierarchy is therefore:
\begin{align}
\epsilon_{\rm infl} &\leftrightarrow \epslamph \quad
  (\text{primary: rate of landscape descent}), \\
\eta_{\rm infl} &\leftrightarrow \etalamph \quad
  (\text{derived: rate of change of descent rate}).
\end{align}
In inflation, $\epsilon$ determines the tensor-to-scalar ratio at leading order via
$r = 16\epsilon$, while $\eta$ controls the tilt via $n_s - 1 \approx -2\epsilon -
\eta$. In RSVP, the conjecture~\eqref{eq:conjecture} assigns to $\epslamph$ the
analogous primary role in determining $\rRSVP$, with $\etalamph$ entering as a
correction. The key difference from inflation is that $\epslamph$ and $\etalamph$
are defined on the field-space manifold $\calF$ via the admissibility volume
functional, without reference to any potential. They characterize the geometry of
the admissibility landscape rather than the curvature of an inflaton potential.

\subsection{Slow Relaxation Condition}

We say that the plenum is in a \emph{slow relaxation} regime when
$|\etalamph| \ll 1$. In this regime, the admissibility gradient is nearly constant,
and the background trajectory in $\calF$ is approximately a geodesic of the
field-space metric $G_{AB}$ deformed by the admissibility gradient. This condition
is the prerequisite for perturbation theory around the lamphrodyne background to be
controlled, in the same way that slow roll is the prerequisite for inflationary
perturbation theory.


%------------------------------------------------------------
\subsection{A Prototype RSVP Action}
\label{sec:action}
%------------------------------------------------------------

A complete cosmological theory requires a dynamical principle from which both the
background evolution and the perturbation equations can be derived. The present
paper does not claim that the final RSVP action has been uniquely identified.
Nevertheless, it is useful to exhibit a minimal prototype action capturing the
essential field content and demonstrating that the formal program is well posed as
an effective field theory.

We consider the provisional action
\begin{equation}
S_{\rm RSVP} = \int d^4x\,\sqrt{-g}\left[
\frac{M_{\rm Pl}^2}{2}R
- \frac{1}{2}\nabla_\mu\Phi\,\nabla^\mu\Phi - V(\Phi)
- \frac{1}{2}\nabla_\mu S\,\nabla^\mu S - W(S)
- \frac{1}{4}F_{\mu\nu}F^{\mu\nu}
+ \xi_1 S\,F_{\mu\nu}F^{\mu\nu}
+ \xi_2\,\nabla_\mu\Phi\,\nabla^\mu S
\right],
\label{eq:action}
\end{equation}
where $F_{\mu\nu} = \nabla_\mu v_\nu - \nabla_\nu v_\mu$ is the field-strength
tensor of the vector field $v_\mu$. The scalar $\Phi$ represents constraint
density, $v_\mu$ represents directed entropy flow, and $S$ represents admissibility
entropy. The interaction terms proportional to $\xi_1$ and $\xi_2$ provide the
lowest-order couplings between the sectors while preserving locality and general
covariance. The Einstein-Hilbert term is included to establish the connection
between RSVP dynamics and the emergent spacetime metric; in the deep lamphrodyne
epoch, the curvature terms are subdominant and the dynamics are dominated by the
plenum sector.

The purpose of this action is not to establish a unique realization of RSVP but to
demonstrate that the framework admits a conventional effective field theory
representation from which perturbative calculations can in principle proceed. We
note a conceptual tension that this prototype does not fully resolve: the paper
claims that global metric expansion is not the smoothing mechanism, yet the action
includes a standard Einstein-Hilbert term $\frac{M_{\rm Pl}^2}{2}R$. Two
resolutions are possible. In the phenomenological embedding (Option B), the metric
is present and dynamical, but the inflationary phase is replaced by lamphrodyne
relaxation as the source of primordial perturbations---the redshift mechanism
is then a subdominant correction from the entropy gradient coupling, not a
replacement for expansion. In the full emergent-metric program (Option A), the
Einstein-Hilbert term is absent and the effective metric is derived from plenum
correlators; this is technically demanding and left for future work. The present
paper adopts the phenomenological embedding for all perturbative calculations,
and all such calculations should be read as estimates within that embedding
rather than consequences of the full emergent-metric claim. The
field-space metric $G_{AB}$ introduced in Section~\ref{sec:fieldspace} is derived
from the kinetic matrix of action~\eqref{eq:action}: at quadratic order in the
fields, $G_{AB}$ takes the diagonal form consistent with the toy metric
\eqref{eq:toy_metric} with $\alpha_S = 1$ and $\alpha_v$ determined by the vector
kinetic term.

%------------------------------------------------------------
\subsection{Euler--Lagrange Equations}
\label{sec:EL}
%------------------------------------------------------------

Variation of the prototype action~\eqref{eq:action} with respect to $\Phi$, $S$,
and $v_\nu$ yields the field equations
\begin{align}
\Box\Phi - V'(\Phi) + \xi_2\,\Box S &= 0, \label{eq:EL_Phi} \\
\Box S - W'(S) + \xi_2\,\Box\Phi + \xi_1\,F_{\mu\nu}F^{\mu\nu} &= 0,
\label{eq:EL_S} \\
\nabla_\mu\!\left[(1 - 4\xi_1 S)\,F^{\mu\nu}\right] &= 0. \label{eq:EL_v}
\end{align}
These equations describe the coupled evolution of the scalar constraint field,
the entropy field, and the directed entropy-flow sector. The coupling $\xi_2$
mixes the $\Phi$ and $S$ equations; the coupling $\xi_1$ introduces a
curvature-like backreaction of the entropy field on the vector sector.

The homogeneous lamphrodyne background corresponds to $\Phi = \Phi_0(t)$,
$S = S_0(t)$, $v_\mu = 0$, for which $F_{\mu\nu} = 0$ identically and the field
equations reduce to the coupled ordinary differential equations
\begin{align}
\ddot\Phi_0 + 3\Heff\dot\Phi_0 + V'(\Phi_0) - \xi_2(\ddot S_0 + 3\Heff\dot S_0)
&= 0, \\
\ddot S_0 + 3\Heff\dot S_0 + W'(S_0) - \xi_2(\ddot\Phi_0 + 3\Heff\dot\Phi_0)
&= 0,
\end{align}
where we have included a friction term proportional to $\Heff$ as the leading-order
coupling of the plenum dynamics to the background rate of entropy-gradient
traversal. These equations are the starting point for the explicit derivation of
the lamphrodyne background trajectory $(\Phi_0(t), S_0(t))$ that the perturbation
calculation requires.

%------------------------------------------------------------
\subsection{Operational Definition of Admissibility}
\label{sec:admissibility_operational}
%------------------------------------------------------------

The admissibility volume $\Vol(\calA)$ introduced in
Section~\ref{sec:fieldspace} becomes physically meaningful only once membership in
$\calA$ is specified by a constitutive condition. For the prototype
theory~\eqref{eq:action}, we define admissible configurations as those satisfying
\begin{equation}
C[\Phi, v, S] \leq 0, \qquad C \equiv \omega^2 - \omega_c^2,
\end{equation}
where $\omega = |\nabla \times \mathbf{v}|$ is the local vorticity magnitude of the
vector field and $\omega_c > 0$ is a critical vorticity threshold. The admissible
region is therefore
\begin{equation}
\calA = \bigl\{(\Phi, \mathbf{v}, S) : |\nabla \times \mathbf{v}| \leq \omega_c
\bigr\}.
\label{eq:admissible_set}
\end{equation}
Configurations with vorticity exceeding $\omega_c$ are interpreted as
thermodynamically unresolved: they carry unresolved rotational structure that the
lamphrodyne dynamics have not yet dissipated. The admissibility volume is then
\begin{equation}
\Vol(\calA) = \int_{C \leq 0} \sqrt{\det G}\; d\Phi\, dS\, d^3v,
\end{equation}
and its decrease in time is controlled by the dissipation of vorticity through the
vector field equation~\eqref{eq:EL_v}. The irreversibility condition
$d\Vol(\calA_{\rm unresolved})/dt < 0$ is therefore not imposed externally but
follows dynamically from equation~\eqref{eq:EL_v}: as vorticity is dissipated, the
admissible region expands and the unresolved complement shrinks.

%------------------------------------------------------------
\subsection{Admissibility Dissipation and Total Loss Budget}
\label{sec:dissipation}
%------------------------------------------------------------

Define the admissibility dissipation rate
\begin{equation}
\Gamma_A = -\frac{d}{dt}\Vol(\calA_{\rm unresolved}) > 0
\end{equation}
during the lamphrodyne epoch. The admissibility evolution equation is
\begin{equation}
\frac{d}{dt}\Vol(\calA_{\rm unresolved}) = -\Gamma_A,
\end{equation}
integrating to
\begin{equation}
\Vol(\calA_{\rm unresolved})(t) = \Vol_0 - \int_0^t \Gamma_A(\tau)\,d\tau.
\end{equation}
By the Monotonic Admissibility Loss Proposition, the integral converges as
$t \to \infty$. Define the total admissibility loss budget
\begin{equation}
\Delta_A = \int_0^\infty \Gamma_A(t)\,dt < \infty.
\label{eq:loss_budget}
\end{equation}
The quantity $\Delta_A$ is a finite, theory-dependent constant that characterizes
the total extent of relaxation undergone by the plenum during the lamphrodyne epoch.
It plays a role analogous to a total entropy-production budget: just as thermodynamic
processes are bounded by the total available free energy, lamphrodyne relaxation is
bounded by $\Delta_A$. A large $\Delta_A$ corresponds to a plenum that undergoes
extensive smoothing; a small $\Delta_A$ corresponds to a nearly pre-resolved initial
state. The connection between $\Delta_A$ and the observed amplitude of primordial
perturbations $A_s$ is a prediction target of the RSVP perturbation program.

%------------------------------------------------------------
\subsection{Admissibility Horizons and Causal Coherence}
\label{sec:horizons}
%------------------------------------------------------------

A central objection to any non-inflationary smoothing mechanism concerns causal
coherence: if regions were never in causal contact, what physically enforces their
approach to a common equilibrium? Inflation answers this by altering causal
structure via accelerated expansion; QQG inherits the same answer. RSVP requires a
different account.

We define an admissibility propagation cone $\mathcal{C}_A$ by the condition
\begin{equation}
ds_A^2 = -c_A^2\,dt^2 + d\mathbf{x}^2 = 0,
\label{eq:admissibility_cone}
\end{equation}
where $c_A$ is the characteristic speed at which admissibility information
propagates through the plenum. The corresponding admissibility horizon is
\begin{equation}
d_A(t) = \int_0^t c_A(\tau)\,d\tau.
\label{eq:admissibility_horizon}
\end{equation}
In the prototype action~\eqref{eq:action}, $c_A$ is determined by the sound speeds
of the coupled $(\Phi, S, v)$ system. If the entropy-field sound speed $c_S$ and
vorticity propagation speed $c_v$ are both superluminal with respect to the
effective metric speed during the deep lamphrodyne epoch---which is possible if the
plenum dynamics precede the emergence of the standard metric---then $d_A(t)$ at
early times exceeds the effective particle horizon of the emergent spacetime
description.

This provides a concrete mechanism by which large-scale coherence may be established
in the plenum before the effective metric description is valid. Regions that appear
causally disconnected in the FLRW description were, in the plenum picture,
connected by admissibility propagation during the epoch when $c_A > c_{\rm eff}$.
This is the RSVP analogue of the inflationary resolution of the horizon problem,
and it makes a testable prediction: the coherence length of primordial perturbations
in RSVP is set by $d_A$ rather than the comoving Hubble radius at horizon crossing,
which implies a distinct dependence of the perturbation spectrum on scale that
differs from the inflationary prediction at very large scales ($k \to 0$).

%------------------------------------------------------------
\section{Perturbations of the Relaxing Plenum}
\label{sec:perturbations}
%------------------------------------------------------------

\subsection{Linearization}

We expand the RSVP fields around the homogeneous background:
\begin{equation}
\Phi = \Phi_0(t) + \delta\Phi(t, \mathbf{x}), \quad
\vv = \delta\vv(t, \mathbf{x}), \quad
S = S_0(t) + \delta S(t, \mathbf{x}).
\end{equation}
The linearized field equations are obtained by varying the RSVP action to first
order in the perturbations. In Fourier space, the perturbations decompose into
scalar, vector, and tensor sectors under the spatial rotation group $SO(3)$.

\subsection{Sector Classification}

The scalar sector contains perturbations in $\delta\Phi$ and $\delta S$, together
with the longitudinal part of $\delta\vv$. These mix under the linearized equations
and produce the scalar power spectrum $P_s(k)$, which must match the observed CMB
anisotropy amplitude $A_s \approx 2.1 \times 10^{-9}$.

The vector sector contains the transverse part of $\delta\vv$. In standard
cosmology, vector perturbations decay during radiation domination and are
observationally irrelevant. In RSVP, where there is no metric expansion, the
transverse vector perturbations do not decay by dilution and may persist to late
times. Their contribution to the CMB anisotropy must be suppressed by the dynamics
of the lamphrodyne epoch, which requires either that $\delta\vv_\perp$ is small
initially or that it decays through entropy production.

The tensor sector contains the transverse-traceless (TT) part of the perturbations.
In the absence of metric fluctuations, tensor perturbations in RSVP arise entirely
from the anisotropic stress of the plenum fields, as we now derive.

%------------------------------------------------------------
\section{Tensor Signatures of Lamphrodyne Relaxation}
\label{sec:tensor}
%------------------------------------------------------------

\subsection{Sources of Tensor Modes}

In standard inflationary cosmology, tensor perturbations---primordial gravitational
waves---arise from quantum fluctuations of the metric during the de Sitter phase.
Their amplitude is set by the Hubble parameter during inflation: $P_h \sim
H_{\rm inf}^2/M_{\rm Pl}^2$. In QQG, the same mechanism operates, with $H_{\rm
inf}$ determined by the RG-improved potential.

In RSVP, no de Sitter phase occurs, so metric quantum fluctuations are not the
source. Tensor perturbations arise instead from the anisotropic stress of the
relaxing plenum. Three sources contribute at quadratic order in the plenum fields.

\subsection{Effective Anisotropic Stress Tensor}

The effective anisotropic stress tensor of the RSVP plenum is
\begin{equation}
\Teff = \alpha\,v_i v_j + \beta\,\partial_i S\,\partial_j S + \gamma\,\partial_i\Phi\,\partial_j\Phi,
\label{eq:Teff}
\end{equation}
where $\alpha$, $\beta$, $\gamma$ are coupling coefficients determined by the RSVP
action, and indices are spatial. This tensor is symmetric and, in general, has both
trace and traceless parts. The trace part contributes to the effective pressure of
the plenum; the transverse-traceless part sources gravitational waves.

The TT projection of~\eqref{eq:Teff} is
\begin{equation}
\Pi_{ij}^{TT} = \Lambda_{ij,kl}\,T^{\rm eff}_{kl},
\end{equation}
where $\Lambda_{ij,kl} = P_{ik}P_{jl} - \frac{1}{2}P_{ij}P_{kl}$ is the TT
projection operator and $P_{ij} = \delta_{ij} - \hat{k}_i\hat{k}_j$ is the
transverse projector for wavevector $\mathbf{k}$.

\subsection{Tensor Power Spectrum}

The tensor power spectrum sourced by the anisotropic stress of the plenum is
\begin{equation}
\Ph(k) = \frac{4}{M_{\rm Pl}^4}\int\frac{d^3q}{(2\pi)^3}\,
G^2_{\rm ret}(k, \tau)\,
\Lambda_{ij,kl}\Lambda_{ij,mn}\,
\langle T^{\rm eff}_{kl}(\mathbf{q})T^{\rm eff}_{mn}(\mathbf{k}-\mathbf{q})\rangle,
\label{eq:Ph}
\end{equation}
where $G_{\rm ret}(k, \tau)$ is the retarded Green function for the tensor mode
equation, and the angle brackets denote the two-point correlator evaluated in the
lamphrodyne background state.

The two-point correlators of the plenum fields in Fourier space are
\begin{align}
\langle\delta v_i(\mathbf{k})\,\delta v_j(\mathbf{k}')\rangle
&= (2\pi)^3\delta^{(3)}(\mathbf{k}+\mathbf{k}')\,P_v(k)\,\delta_{ij}, \\
\langle\delta S(\mathbf{k})\,\delta S(\mathbf{k}')\rangle
&= (2\pi)^3\delta^{(3)}(\mathbf{k}+\mathbf{k}')\,P_S(k), \\
\langle\delta\Phi(\mathbf{k})\,\delta\Phi(\mathbf{k}')\rangle
&= (2\pi)^3\delta^{(3)}(\mathbf{k}+\mathbf{k}')\,P_\Phi(k),
\end{align}
where $P_v(k)$, $P_S(k)$, $P_\Phi(k)$ are the individual power spectra of the
plenum field fluctuations, to be derived from the linearized field equations of the
lamphrodyne epoch.

\subsection{Tensor-to-Scalar Ratio}

The tensor-to-scalar ratio in RSVP is defined as
\begin{equation}
\rRSVP = \frac{\Ph(k_*)}{P_s(k_*)},
\label{eq:rRSVP}
\end{equation}
where $k_* \approx 0.05\,{\rm Mpc}^{-1}$ is the CMB pivot scale. Substituting
the contributions from each plenum field source,
\begin{equation}
\rRSVP = F(\alpha, \beta, \gamma, P_v, P_S, P_\Phi, \etalamph),
\label{eq:rRSVP_general}
\end{equation}
where $F$ is a functional that depends on the plenum coupling coefficients, the
individual field power spectra, and the lamphrodyne relaxation parameter. In the
slow-relaxation limit $|\etalamph| \ll 1$, $\rRSVP$ simplifies and can be expressed
as a power law in $\etalamph$, analogously to the inflationary consistency relation
$r = 16\epsilon$.

The derivation of the full functional $F$ requires the explicit power spectra
$P_v$, $P_S$, $P_\Phi$, which in turn require solving the linearized RSVP field
equations on the lamphrodyne background. This is the primary open calculation
of the present framework; we present the formal structure here and note that
the derivation of $\rRSVP$ as a function of $\etalamph$ is the analog of deriving
$r = 16\epsilon$ in inflation---a result that required the full apparatus of
inflationary perturbation theory and was not available at the time the inflationary
paradigm was first proposed.

\subsection{Tensor Suppression and Spectral Tilt}

The tensor spectral index $n_t$ is defined by $\Ph(k) \propto k^{n_t}$. In
slow-roll inflation $n_t = -r/8 \approx -2\epsilon$, giving a slightly red
spectrum. In QQG the tilt is similarly red. In RSVP, the tensor spectrum is sourced
by plenum anisotropic stress rather than de Sitter vacuum fluctuations, which
generically produces a different tilt.

A key feature of the prototype action is the coupling $\xi_1 S F_{\mu\nu}F^{\mu\nu}$,
which modifies the effective gauge kinetic function of the vector field to
$(1 - 4\xi_1 S)$. This is a kinetic rescaling, not a Proca mass. It does not
generate an $m_v^2 v_\mu v^\mu$ term; producing a genuine vector mass would require
an additional Stückelberg or Higgs sector not present in the current prototype
action. The distinction matters for the tensor suppression mechanism: a kinetic
rescaling by factor $(1 - 4\xi_1 S_0)$ suppresses the vector propagator amplitude
by $|1 - 4\xi_1 S_0|^{-1}$ relative to the unrescaled case, while a Proca mass
suppresses it by $(k^2 + m_v^2)^{-1/2}$. The suppression mechanisms are
qualitatively different.

Within the prototype action as written, the tensor amplitude is suppressed relative
to the inflationary vacuum amplitude by the smallness of the anisotropic stress
correlators. The effective coupling ratio
\begin{equation}
\mathcal{R}_{\rm sup} \equiv \frac{\alpha^2}{|1 - 4\xi_1 S_0|^2} \cdot
\frac{\Heff^2}{M_{\rm Pl}^4}
\label{eq:suppression}
\end{equation}
controls the overall amplitude of $\Ph(k)$ relative to the scalar spectrum. For
$(1 - 4\xi_1 S_0)$ of order unity and $\alpha \ll M_{\rm Pl}$, the tensor-to-scalar
ratio is suppressed relative to the inflationary value $r = 16\epslamph$. The
exact suppression depends on the background value of $S_0$ during the lamphrodyne
epoch, which requires solving the full background equations~\eqref{eq:EL_Phi}
--\eqref{eq:EL_S}.

Regarding the spectral tilt: the $k$-dependence of $\Ph(k)$ is inherited from
the $k$-dependence of the plenum correlators. For a vector field with kinetic
function $(1 - 4\xi_1 S_0)$ varying slowly in time, the vector power spectrum
$P_v(k)$ will be approximately scale-invariant at scales $k \gg \Heff$ and will
turn over at $k \sim \Heff$. The resulting tensor tilt $n_t^{\rm RSVP}$ is
therefore expected to be near zero or slightly positive (blue) at scales
$k \gtrsim \Heff$, in qualitative contrast to inflation's red tilt. A blue tensor
tilt is a structural prediction of sourced tensor mechanisms and would provide
a distinctive observational signature of Class~3 smoothing, independent of the
amplitude.

\subsection{Admissibility Crossing as the Native Freeze-Out Condition}

A more fundamental issue for the perturbation theory underlies both the scalar and
tensor calculations. In inflationary perturbation theory, modes freeze when $k =
aH$: the mode wavelength exceeds the Hubble horizon, and perturbations stop
evolving. In QQG the same mechanism operates because the inflationary phase is
still de Sitter. In RSVP, there is no metric expansion, so Hubble horizon crossing
does not occur in the standard sense. The scalar spectrum derivation in
Section~\ref{sec:scalar_spectrum} used the phenomenological GR embedding precisely
to have access to $k = aH$ as a freeze-out condition, at the cost of importing FRW
and Friedmann dynamics.

The RSVP-native freeze-out condition is admissibility crossing: a mode $k$ freezes
when
\begin{equation}
k \cdot d_A(t) \sim 1,
\label{eq:admissibility_crossing}
\end{equation}
where $d_A(t) = \int_0^t c_A(\tau)\,d\tau$ is the admissibility horizon. Modes
with $k \cdot d_A > 1$ are within the admissibility horizon and continue to
exchange information through plenum dynamics; modes with $k \cdot d_A < 1$ have
exited the admissibility horizon and are effectively frozen. This is the RSVP
analogue of Hubble horizon crossing, with $d_A$ replacing $1/H$.

The spectral index evaluated at admissibility crossing will differ from the
phenomenological estimate~\eqref{eq:ns_conjecture} in precisely the regime where
$d_A$ departs from $1/\Heff$. At large scales ($k \to 0$), if $c_A$ is
approximately constant during the lamphrodyne epoch, then $d_A \propto t$ and the
spectrum acquires a scale dependence controlled by the time evolution of
$\epslamph(t)$. Deriving $n_s$ from admissibility crossing rather than Hubble
crossing is the calculation that would produce a genuinely RSVP-native spectral
index, as opposed to the multifield analogue estimate of
Section~\ref{sec:scalar_spectrum}.

\begin{proposition}[Lamphrodyne Consistency Conjecture]
In the slow-relaxation regime $|\etalamph| \ll 1$, the tensor-to-scalar ratio
admits an expansion
\begin{equation}
\rRSVP = C\,\epslamph^{\,p} + \mathcal{O}(\etalamph),
\label{eq:conjecture}
\end{equation}
where the coefficient $C$ and exponent $p$ depend on the dominant plenum source
contributing to $\Teff$. Specifically, if vector vorticity dominates, $p$ is
determined by the scaling of $P_v(k)$ with $\epslamph$; if entropy shear dominates,
$p$ reflects the scaling of $P_S(k)$. The slow-relaxation approximation
$|\etalamph| \ll 1$ is the condition under which the $\mathcal{O}(\etalamph)$
correction is negligible, exactly as $|\eta_{\rm infl}| \ll 1$ is required for the
inflationary consistency relation $r = 16\epsilon$ to hold.
\end{proposition}

This conjecture is not proven here; it defines the research program. The derivation
of the exponent $p$ and coefficient $C$ from explicit plenum power spectra is the
calculation that would transform the present framework paper into a predictive
cosmological model. We state the conjecture explicitly so that its resolution can
serve as a milestone for the RSVP program, analogous to the derivation of
$r = 16\epsilon$ in inflationary perturbation theory.

%------------------------------------------------------------
\section{The Empirical Fork}
\label{sec:fork}
%------------------------------------------------------------

\subsection{The QQG Tensor Floor}

The structural prediction of QQG is equation~\eqref{eq:qqg_floor}: the
tensor-to-scalar ratio satisfies $\rQQG \gtrsim 10^{-2}$ as a consequence of the
perturbativity bound $\lambda_{tH} \lesssim 1$. This bound is not a
phenomenological choice but an internal consistency condition: if $\lambda_{tH}$
exceeds unity, the 1-loop approximation that underlies the entire QQG inflationary
scenario breaks down. The tensor floor is therefore a necessary consequence of the
theory's validity conditions.

The next generation of CMB polarization experiments---LiteBIRD, CMB-S4, and the
Simons Observatory---will reach sensitivities of $\sigma(r) \sim 10^{-3}$, well
below the QQG floor. A robust upper bound at $r < 10^{-2}$ would strongly challenge the simplest QQG
realization and require modification of its perturbative assumptions.

\subsection{RSVP Predictions: Three Cases}

The RSVP prediction for $\rRSVP$ depends on the specific dynamics of the
lamphrodyne epoch, which are not yet fully constrained. Three qualitative cases
are distinguished.

\textit{Case A: Suppressed tensor spectrum.} If the plenum field correlators
$P_v$, $P_S$, $P_\Phi$ are strongly suppressed relative to the scalar power
spectrum during the lamphrodyne epoch, then $\rRSVP \ll 10^{-2}$. This case is
consistent with all current B-mode non-detections and would be favored if
future experiments find no primordial tensor signal.

\textit{Case B: Inflation-like tensor spectrum.} If the anisotropic stress of the
relaxing plenum generates tensor modes at an amplitude comparable to the QQG floor,
then $\rRSVP \sim 10^{-2}$. In this case, RSVP must account for why relaxation
produces a tensor residue that mimics de Sitter quantum fluctuations in amplitude.

\textit{Case C: Distinct spectral shape.} The tensor spectrum $\Ph(k)$ may have a
spectral tilt $n_t^{\rm RSVP}$ that differs from the inflationary prediction
$n_t = -2\epsilon$. Even if the amplitude $\rRSVP$ falls within the observationally
accessible range, a blue or unusually tilted tensor spectrum would be a distinctive
signature of lamphrodyne relaxation rather than de Sitter inflation.

\subsection{Falsification Matrix}

Table~\ref{tab:fork} summarizes the observational consequences for QQG and RSVP
across the principal experimental outcomes.

\begin{table}[h]
\centering
\caption{Empirical fork: observational outcomes and their implications for
Class~2 (QQG) and Class~3 (RSVP) smoothing mechanisms.}
\label{tab:fork}
\begin{tabular}{>{\raggedright}p{5cm} p{4cm} p{4cm}}
\toprule
\textbf{Observation} & \textbf{QQG} & \textbf{RSVP} \\
\midrule
$r > 0.01$ & Survives; constrains $\lambda_{tH}$ & Must explain tensor residue
from plenum stress \\[4pt]
$r \approx 0.01$ & Survives at parameter boundary & Constrained; Case~B viable \\[4pt]
$r < 0.01$ & Falsified by internal bound & Survives; Case~A or~C favored \\[4pt]
Unusual tilt $n_t > 0$ & Requires significant model extension & Natural if plenum
correlators are blue \\[4pt]
Tilt consistent with $n_t = -r/8$ & Standard consistency; favors inflation &
Coincidental; requires explanation \\
\bottomrule
\end{tabular}
\end{table}

%------------------------------------------------------------
\section{Ontological Economy and the Large-\texorpdfstring{$N$}{N} Question}
\label{sec:economy}
%------------------------------------------------------------

\subsection{Why QQG Requires Large $N$}

The QQG inflationary scenario requires $N \sim 10^5$--$10^6$ matter fields to
produce a spectral index $n_s$ compatible with CMB observations. This requirement
is not optional: in the vacuum ($N = 0$) case, the beta-function trajectories do
not generate a viable inflationary potential \citep{Liu2026}. The matter fields
enter only through their vacuum fluctuation contributions to the one-loop beta
functions; they need not be excited, and their physical content beyond loop
corrections is left unspecified.

The 't~Hooft-like coupling $\lambda_{tH} = \lambda_0 N/(4\pi)^2 \lesssim 1$ provides
the perturbativity bound. As $N$ increases, $\lambda_0$ must decrease
proportionally to keep $\lambda_{tH}$ under control. The large-$N$ matter content
is therefore not a prediction of the theory but a parameter space requirement: it
is needed to make the beta functions produce a viable slow-roll trajectory.

\subsection{The Empirical Status of the Large-$N$ Sector}

No independent observational consequence of the large-$N$ sector has yet been
identified. The matter fields sit in their vacua throughout the inflationary epoch;
they do not contribute to reheating, do not produce distinctive non-Gaussianities,
and do not generate spectral features that would independently confirm their
presence. Their role is to modify the RG flow of the gravitational couplings $\xi$
and $\lambda$ so that a viable slow-roll trajectory is obtained.

This is a meaningful distinction. The inflaton fields of Class~1 smoothing are
empirically unjustified in their potential shape, but they do make independent
predictions about reheating dynamics, curvaton mechanisms, and spectral features
that can in principle be tested. Whether the large-$N$ spectator sector of QQG
will ultimately yield analogous independent predictions---through, for example,
its role in the strong-coupling regime or in the particle spectrum of the emergent
radiation era---is an open question. As of the present paper, the sector is
mathematically motivated by the requirements of the beta-function structure, while
its empirical motivation remains to be established.

\subsection{Relaxation Cosmology as an Alternative}

The three RSVP plenum fields $\Phi$, $\vv$, $S$ are not introduced to fix
observational parameters. They are required by the theoretical structure of the
admissibility and constraint-closure program: $\Phi$ encodes constraint
satisfaction, $\vv$ encodes directed entropy flux, and $S$ encodes
irreversibility. Each field carries independent theoretical motivation; none is
introduced solely to modify a running coupling.

The ontological economy argument is not a simple appeal to parsimony or field
count. A more precise formulation is the following: the QQG large-$N$ sector
participates in the cosmological dynamics exclusively through renormalization-group
effects on the gravitational couplings. Its fields sit in their vacua; they do not
appear in the background equations, the reheating dynamics, or the perturbation
spectrum except indirectly through beta-function corrections. The RSVP plenum fields
$\Phi$, $\vv$, $S$, by contrast, participate directly in the dynamics being
modeled: they define the background state, source the perturbations, generate the
effective stress tensor, and carry the entropy gradient that produces the observed
redshift. The distinction is not between many fields and few fields, but between
fields that are dynamically active in the mechanism being proposed and fields that
function as renormalization-group auxiliaries. Both frameworks introduce speculative
field content; the claim here is only that RSVP's field content is dynamically
necessary to the mechanism in a way that the QQG large-$N$ sector presently is not.

%------------------------------------------------------------
\section{Discussion}
\label{sec:discussion}
%------------------------------------------------------------

\subsection{Inflation After the Inflaton}

The QQG result of \citet{Liu2026} marks a genuine advance in theoretical
cosmology. It demonstrates that the inflationary phase is not necessarily tied to
the existence of an auxiliary scalar field with a tuned potential, but can instead
emerge from the UV-complete dynamics of quantum gravity itself. This is a
conceptually important result independent of whether QQG inflation ultimately
proves correct: it establishes that the inflaton paradigm is not the only way to
generate a quasi-de Sitter phase, and that the physical source of inflationary
energy can be sought within the gravitational sector rather than appended to it.

\subsection{What Remains to Be Explained}

The QQG scenario does not address the question of whether inflationary expansion
is necessary, only the question of its source. The three fundamental requirements
of early-universe smoothing---isotropy, homogeneity, and scale-invariant
perturbations---are explained within QQG by the same mechanism as in Class~1
inflation: a de Sitter phase stretches modes beyond the Hubble horizon and
establishes causal coherence before last scattering. QQG changes the energy
source of this phase; it does not change its physical mechanism.

RSVP proposes that the mechanism itself can be replaced. This is a stronger claim,
and it carries a heavier formal burden. The present paper takes a step toward
discharging that burden by providing the formal structure of the lamphrodyne
background, the definition of $\etalamph$, and the tensor spectrum formula
$\rRSVP$. The full derivation of $\rRSVP$ as a function of $\etalamph$---the
RSVP analogue of $r = 16\epsilon$---remains the primary open problem of the
framework.

\subsection{Toward an RSVP Consistency Relation}

The inflation literature is organized around the consistency relation $r =
16\epsilon$. QQG organizes its predictions around the structural bound $\rQQG
\gtrsim 10^{-2}$. RSVP is proposing the conjecture~\eqref{eq:conjecture}:
\begin{equation}
\rRSVP = C\,\epslamph^{\,p} + \mathcal{O}(\etalamph).
\end{equation}
These three relations define the comparison between smoothing classes at the level
of a single observable:
\begin{align}
\text{Class 1 (inflation):} &\quad r = 16\epsilon_{\rm infl}, \\
\text{Class 2 (QQG):} &\quad r \gtrsim 10^{-2}, \\
\text{Class 3 (RSVP):} &\quad r = C\,\epslamph^{\,p} + \mathcal{O}(\etalamph).
\end{align}
The inflation relation is derived; the QQG bound is structural; the RSVP conjecture
is programmatic. Deriving the exponent $p$ and coefficient $C$ from explicit plenum
power spectra $P_v(k)$, $P_S(k)$, $P_\Phi(k)$ is the calculation that promotes
the RSVP entry from programmatic to derived. At that point the comparison becomes
genuinely quantitative rather than structural, and the three-class taxonomy becomes
a three-model competition with a shared observable target.

The sequence of steps required is clear from the framework developed here: specify
the RSVP action and kinetic terms to fix $G_{AB}$; solve the linearized field
equations on the lamphrodyne background to derive $P_v(k)$, $P_S(k)$, $P_\Phi(k)$;
project $\Teff$ onto the TT sector; evaluate $\Ph(k)$; and extract $p$ and $C$ in
the slow-relaxation limit. This is a well-posed calculation within the framework
established here.

\subsection{Multifield Analogy and Preliminary Spectral Estimate}
\label{sec:scalar_spectrum}

The scalar power spectrum is observationally primary. Any candidate cosmological
model must reproduce the observed amplitude $A_s \approx 2.1 \times 10^{-9}$ and
spectral index $n_s \approx 0.965$ \citep{Planck2020} before tensor predictions
carry weight. We are at an earlier stage for scalar perturbations than for tensor
perturbations, and we flag this explicitly. What follows is a heuristic estimate
based on embedding RSVP in a phenomenological GR context; it is not a derived
result from the RSVP-native perturbation theory, which requires admissibility
horizons rather than Hubble horizons as the fundamental freeze-out condition. We
present it as a multifield analogy to indicate the order of magnitude of what a
complete calculation might yield.

The prototype action~\eqref{eq:action}, restricted to its scalar sector
$(\Phi, S)$ on a homogeneous FRW background, describes a two-field system with
non-canonical kinetic mixing. The field-space metric in the $(\Phi, S)$ subspace
is read off from the kinetic terms:
\begin{equation}
G_{AB} = \begin{pmatrix} 1 & -\xi_2 \\ -\xi_2 & 1 \end{pmatrix},
\qquad
G^{AB} = \frac{1}{1-\xi_2^2}\begin{pmatrix} 1 & \xi_2 \\ \xi_2 & 1 \end{pmatrix}.
\end{equation}
In the phenomenological GR embedding---where we treat the background as FRW and use
$H$ for the metric Hubble rate---the slow-roll parameter $\epsilon_{\rm lamph}$,
interpreted via the field-space gradient of the effective potential $V_{\rm eff} =
V(\Phi) + W(S)$, takes the form
\begin{equation}
\epsilon_{\rm lamph} \approx \frac{M_{\rm Pl}^2}{2}
\frac{(V')^2 + (W')^2 + 2\xi_2 V'W'}{(1-\xi_2^2)(V+W)^2}.
\label{eq:eps_operational}
\end{equation}
This connects the admissibility-gradient definition of $\epslamph$ to an operational
expression in terms of the prototype potentials, under the assumption that
$\nabla_{\calF}\log\Vol(\calA)$ is proportional to $\nabla_{\calF}V_{\rm eff}$ for
the vorticity-based admissibility constraint~\eqref{eq:admissible_set}. The
proportionality holds when the admissibility boundary is level sets of $V_{\rm eff}$
in field space, which is an assumption whose justification requires the full
admissibility dynamics.

Using the standard multi-field perturbation theory of \citet{Gordon2001} for a
two-scalar system with field-space metric $G_{AB}$, the scalar spectral index at
leading order in slow roll is
\begin{equation}
n_s - 1 = -a\,\epslamph - b\,\etalamph + \mathcal{O}(\epslamph^2),
\label{eq:ns_conjecture}
\end{equation}
where the coefficients $a$ and $b$ depend on the trajectory curvature in field
space, the isocurvature-to-adiabatic transfer, and the mixing parameter $\xi_2$.
In the single-field limit $\xi_2 \to 0$, the standard result $a = 6$, $b = -2$
is recovered. For nonzero $\xi_2$, the coefficients shift by $\mathcal{O}(\xi_2^2)$
corrections.

We emphasize that equation~\eqref{eq:ns_conjecture} is derived within the
phenomenological GR embedding, not from the RSVP-native perturbation theory. In
the GR embedding, horizon crossing at $k = aH$ determines the freeze-out scale,
and the Mukhanov--Sasaki formalism applies directly. In the full RSVP framework,
freeze-out instead occurs at admissibility crossing $k\,d_A(t) \sim 1$, where
$d_A$ is the admissibility horizon defined in Section~\ref{sec:horizons}. The
RSVP-native spectral index will therefore differ from~\eqref{eq:ns_conjecture}
in its scale dependence at $k \to 0$, where the admissibility horizon structure
modifies the spectrum relative to the inflationary prediction. Deriving this
correction is the principal remaining calculation of the scalar perturbation
program.

Define the scalar power spectrum
\begin{equation}
P_s(k) = \langle \zeta(\mathbf{k})\,\zeta(-\mathbf{k})\rangle,
\end{equation}
where $\zeta$ is the gauge-invariant curvature perturbation. In the
phenomenological embedding, $P_s(k_*) \approx H^2/(8\pi^2 M_{\rm Pl}^2
\epslamph)$ at the pivot scale $k_*$, matching the inflationary formula with
$\epsilon_{\rm infl}$ replaced by $\epslamph$. Setting this equal to $A_s$ gives
a constraint on $\epslamph$ as a function of $H$ during the lamphrodyne epoch.
This is the scalar normalization condition that the background solution must
satisfy, and it connects $\epslamph$ to the observed CMB amplitude independently
of the tensor sector.

\subsection{Falsifiability of RSVP}

It is important to state explicitly what would falsify the RSVP framework, not
merely challenge it. Class~3 smoothing would be effectively refuted if all
observable tensor properties satisfy inflationary consistency relations to high
precision: specifically, if the tensor spectral index is measured to satisfy
$n_t = -r/8$ with high significance, this would confirm de Sitter vacuum
fluctuations as the tensor source and leave no room for a sourced alternative.
RSVP would also be challenged if no admissibility-based relaxation model can
simultaneously reproduce the observed scalar power spectrum amplitude $A_s$ and
spectral index $n_s$ while maintaining the irreversibility condition
$d\Vol(\calA_{\rm unresolved})/dt < 0$ throughout the lamphrodyne epoch---a failure
of internal consistency analogous to QQG's strong-coupling bound. Finally, if the
tensor spectrum is shown to be generated prior to any thermodynamic epoch, in a
regime where entropy production is observationally excluded, the lamphrodyne
mechanism loses its physical grounding. These are not remote possibilities; they
are concrete conditions that the derivation of $\rRSVP$ must be tested against.

\subsection{Future Observations}

The observational landscape over the next decade will provide significant
discriminating power. LiteBIRD, scheduled for launch in the early 2030s, targets
a sensitivity of $\sigma(r) \sim 10^{-3}$ and will either detect primordial tensor
modes above the QQG floor or push below it. CMB-S4 will provide complementary
constraints on $n_s$, non-Gaussianity, and CMB spectral distortions. The Simons
Observatory is currently operational and will provide intermediate constraints.

For RSVP, the key observational target is not a single value of $r$ but the
tensor spectral tilt $n_t$. If primordial tensor modes are detected, their spectral
shape will distinguish de Sitter quantum fluctuations (slightly red, $n_t \approx
-r/8$) from sourced tensor modes of the kind RSVP predicts. Specifically, if the
tensor spectrum is blue ($n_t > 0$) at the pivot scale, this cannot be explained
within any de Sitter inflation framework and would constitute strong evidence for
a non-inflationary source of tensor modes, which lamphrodyne relaxation is
designed to provide.

%------------------------------------------------------------
\section{Conclusion}
\label{sec:conclusion}
%------------------------------------------------------------

We have proposed a classification of early-universe smoothing mechanisms into
three structurally distinct classes. Class~1---inflaton smoothing---uses an
auxiliary scalar field with tuned potential to drive metric expansion. Class~2---
geometric smoothing, exemplified by the QQG result of \citet{Liu2026}---uses
quantum-corrected gravitational dynamics to generate the same expansion without a
separate inflaton. Class~3---lamphrodyne relaxation, proposed within the RSVP
framework---uses admissibility-reducing dynamics in a coupled scalar-vector-entropy
plenum to achieve cosmological homogenization without global metric expansion.

We have formalized Class~3 by defining the lamphrodyne background state $(\Phi_0,
\vv_0, S_0)$, imposing the irreversibility condition
$d\Vol(\calA_{\rm unresolved})/dt < 0$, introducing the field-space admissibility
metric on $\calF$, and defining the lamphrodyne relaxation parameter $\etalamph$
from admissibility geometry. We have constructed the effective anisotropic stress
tensor of the relaxing plenum and constructed the formal structure of the tensor power spectrum
$\Ph(k)$ and tensor-to-scalar ratio $\rRSVP$. We have compared these with the
QQG tensor floor $\rQQG \gtrsim 10^{-2}$ and identified the empirical fork that
future B-mode experiments will resolve.

The significance of Quadratic Quantum Gravity may not be that it explains
inflation. Its deeper significance is that it demonstrates that cosmological
smoothing can emerge from the intrinsic dynamics of a fundamental theory. Once
this possibility is admitted, inflationary expansion becomes one member of a
broader class of smoothing mechanisms. RSVP explores a different member of this
class: relaxation without expansion. Future measurements of primordial tensor
modes---their amplitude, spectral tilt, and scale dependence---may determine which
description better captures the universe's earliest history.

The primary contribution of this paper is the demonstration that Class~3 smoothing
is formally tractable. The lamphrodyne relaxation parameter $\etalamph$ plays the
same structural role in RSVP that the slow-roll parameters play in inflation: it
characterizes the regime in which the background is stable, perturbation theory is
valid, and observables can be computed. The derivation of $\rRSVP$ as an explicit
function of $\etalamph$---the RSVP consistency relation---is the next step, and it
is a step that can be taken within the framework presented here.

%------------------------------------------------------------
\bibliographystyle{plainnat}
\begin{thebibliography}{99}

\bibitem[Albrecht and Steinhardt(1982)]{Albrecht1982}
Albrecht, A. and Steinhardt, P.~J. (1982).
Cosmology for grand unified theories with radiatively induced symmetry breaking.
\textit{Physical Review Letters}, 48(17):1220--1223.

\bibitem[Avramidi and Barvinsky(1985)]{Avramidi1985}
Avramidi, I.~G. and Barvinsky, A.~O. (1985).
Asymptotic freedom in higher-derivative quantum gravity.
\textit{Physics Letters B}, 159(4):269--274.

\bibitem[BICEP/Keck Collaboration(2021)]{BICEPKeck2021}
Ade, P.~A.~R. et al. (2021).
Improved constraints on primordial gravitational waves using Planck, WMAP, and BICEP/Keck observations through the 2018 observing season.
\textit{Physical Review Letters}, 127(15):151301.

\bibitem[Buccio et al.(2024)]{Buccio2024}
Buccio, D., Donoghue, J.~F., Menezes, G., and Percacci, R. (2024).
Physical running of couplings in quadratic gravity.
\textit{Physical Review Letters}, 133(2):021604.

\bibitem[Capper and Duff(1974)]{Capper1974}
Capper, D.~M. and Duff, M.~J. (1974).
Trace anomalies in dimensional regularization.
\textit{Nuovo Cimento A}, 23:173.

\bibitem[Fradkin and Tseytlin(1982)]{Fradkin1982}
Fradkin, E.~S. and Tseytlin, A.~A. (1982).
Renormalizable asymptotically free quantum theory of gravity.
\textit{Nuclear Physics B}, 201(3):469--491.

\bibitem[Guth(1981)]{Guth1981}
Guth, A.~H. (1981).
Inflationary universe: A possible solution to the horizon and flatness problems.
\textit{Physical Review D}, 23(2):347--356.

\bibitem[Holdom and Ren(2016)]{Holdom2016}
Holdom, B. and Ren, J. (2016).
QCD analogy for quantum gravity.
\textit{Physical Review D}, 93(12):124030.

\bibitem[Kofman et al.(1997)]{Kofman1997}
Kofman, L., Linde, A., and Starobinsky, A.~A. (1997).
Towards the theory of reheating after inflation.
\textit{Physical Review D}, 56(6):3258--3295.

\bibitem[Linde(1982)]{Linde1982}
Linde, A.~D. (1982).
A new inflationary universe scenario.
\textit{Physics Letters B}, 108(6):389--393.

\bibitem[Liu et al.(2026)]{Liu2026}
Liu, R., Quintin, J., and Afshordi, N. (2026).
Ultraviolet completion of the big bang in quadratic gravity.
\textit{Physical Review Letters}, 136(11):111501.

\bibitem[Louis et al.(2025)]{Louis2025}
Louis, T. et al. (ACT Collaboration) (2025).
The Atacama Cosmology Telescope: DR6 power spectra, likelihoods and cosmological parameters.
\textit{Journal of Cosmology and Astroparticle Physics}, 11:062.

\bibitem[Martin et al.(2014)]{Martin2014}
Martin, J., Ringeval, C., and Vennin, V. (2014).
Encyclop{\ae}dia inflationaris.
\textit{Physics of the Dark Universe}, 5--6:75--235.

\bibitem[Planck Collaboration(2020)]{Planck2020}
Akrami, Y. et al. (Planck Collaboration) (2020).
Planck 2018 results X: Constraints on inflation.
\textit{Astronomy and Astrophysics}, 641:A10.

\bibitem[Starobinsky(1980)]{Starobinsky1980}
Starobinsky, A.~A. (1980).
A new type of isotropic cosmological models without singularity.
\textit{Physics Letters B}, 91(1):99--102.

\bibitem[Gordon et al.(2001)]{Gordon2001}
Gordon, C., Wands, D., Bassett, B.~A., and Maartens, R. (2001).
Adiabatic and entropy perturbations from inflation.
\textit{Physical Review D}, 63(2):023511.

\end{thebibliography}

\end{document}
